\begin{ex}
Dare la macchina di Turing che complementa una sequenza binaria e si riposiziona all'inizio della parola.
\end{ex}

\begin{sol}
La macchina di turing è la seguente
\[
\begin{tikzpicture}[
    scale=0.8, transform shape,
    >=stealth', shorten >=1pt, auto,
    every state/.style={circle, draw, fill=white, minimum size=0.8cm, inner sep=1pt}
]
% Stati
\node[state]            (q0) at (0,0)  {$q_0$};
\node[state]            (q1) at (5,0)  {$q_1$};
\node[state, accepting] (qF) at (10,0) {$q_F$};

% Freccia d'ingresso
\draw[->] ([shift={(-0.7,0.7)}]q0.center) -- (q0);

% Cappi
\draw[->] (q0) to[loop above, looseness=8] node {$0/1,+1$} (q0);
\draw[->] (q0) to[loop below, looseness=8] node {$1/0,+1$} (q0);
\draw[->] (q1) to[loop above, looseness=8] node {$0/0,-1$} (q1);
\draw[->] (q1) to[loop below, looseness=8] node {$1/1,-1$} (q1);

% Transizioni
\draw[->] (q0) -- node {$\square/\square,-1$} (q1);
\draw[->] (q1) -- node {$\square/\square,+1$} (qF);
\end{tikzpicture}
\]
forniamo anche il codice in Wolfram Matematica
\begin{walfram}
turing = {{
 {q0, 0} -> {q0, 1, +1},
 {q0, 1} -> {q0, 0, +1},
 {q0, B} -> {q1, B, -1},
 {q1, 0} -> {q1, 0, -1},
 {q1, 1} -> {q1, 1, -1},
 {q1, B} -> {qF, B, +1}
 }, q0, {qF}}
\end{walfram}
\end{sol}

\begin{ex}
Come l'esercizio precedente, con l'unica differenza che la computazione termina solo con la parola vuota e in caso contrario, dopo essere tornata alla posizione iniziale, ripete il ciclo.
\end{ex}
\begin{sol}
La macchina di turing è la seguente
\[
\begin{tikzpicture}[
    scale=0.8, transform shape,
    >=stealth', shorten >=1pt, auto,
    every state/.style={circle, draw, fill=white, minimum size=0.8cm, inner sep=1pt}
]
% Stati della riga principale
\node[state] (q0) at (0,0)  {$q_0$};
\node[state] (q1) at (5,0)  {$q_1$};
\node[state] (qR) at (10,0) {$q_R$};

% Stato finale
\node[state, accepting] (qF) at (0,-3.5) {$q_F$};

% Freccia d'ingresso
\draw[->] ([shift={(-0.7,0.7)}]q0.center) -- (q0);

% Transizioni da q_0 a q_1 (complemento del primo simbolo)
\draw[->] (q0) to[bend left=20]  node {$1/0,+1$} (q1);
\draw[->] (q0) to[bend right=20] node[swap] {$0/1,+1$} (q1);

% Parola vuota: accettazione
\draw[->] (q0) -- node[swap] {$\square/\square,+1$} (qF);

% Cappi
\draw[->] (q1) to[loop above, looseness=8] node {$1/0,+1$} (q1);
\draw[->] (q1) to[loop below, looseness=8] node {$0/1,+1$} (q1);
\draw[->] (qR) to[loop above, looseness=8] node {$0/0,-1$} (qR);
\draw[->] (qR) to[loop below, looseness=8] node {$1/1,-1$} (qR);

% Fine parola: ritorno verso sinistra
\draw[->] (q1) -- node {$\square/\square,-1$} (qR);

% Ritorno a q_0 (ripete il ciclo)
\draw[->] (qR) to[bend left=35] node[below] {$\square/\square,+1$} (q0);
\end{tikzpicture}
\]
forniamo anche il codice in Wolfram Matematica
\begin{walfram}
mu = {{
 {q0,  0} -> {q1, 1, +1},
 {q0,  1} -> {q1, 0, +1},
 {q0, -1} -> {qF, -1, +1},
 {q1,  0} -> {q1, 1, +1},
 {q1,  1} -> {q1, 0, +1},
 {q1, -1} -> {qR, -1, -1},
 {qR,  0} -> {qR, 0, -1},
 {qR,  1} -> {qR, 1, -1},
 {qR, -1} -> {q0, -1, +1}
 }, q0, {qF}}
\end{walfram}
\end{sol}